% ===========================================================================
% Paper 50: Algorithmic Number Theory
% Specialist Mathematics Project
% @author: Michael Fuhrmann
% @date: 2026-03-12
% @build: 137
% ===========================================================================

\documentclass[12pt,a4paper]{amsart}
\usepackage{amsmath,amssymb,amsthm}
\usepackage{mathtools}
\usepackage[utf8]{inputenc}
\usepackage[T1]{fontenc}
\usepackage{hyperref}
\usepackage{enumitem,booktabs,array}
\usepackage{geometry}
\geometry{margin=2.5cm}

\newtheorem{theorem}{Theorem}[section]
\newtheorem{lemma}[theorem]{Lemma}
\newtheorem{corollary}[theorem]{Corollary}
\newtheorem{proposition}[theorem]{Proposition}
\newtheorem{conjecture}[theorem]{Conjecture}
\theoremstyle{definition}
\newtheorem{definition}[theorem]{Definition}
\newtheorem{example}[theorem]{Example}
\newtheorem{remark}[theorem]{Remark}

\author{Michael Fuhrmann}
\address{Specialist Mathematics Project, Build 137}
\email{specialist-maths@localhost}
\date{\today}

\title{Algorithmic Number Theory:\\
Foundations, Complexity, and Cryptographic Applications}

\begin{document}

\begin{abstract}
Algorithmic number theory investigates the computational aspects of integer
arithmetic: primality testing, integer factorization, discrete logarithms, and
the number-theoretic foundations of modern cryptography. This paper surveys
the principal algorithms and complexity results organized around five core
topics: the Euclidean algorithm and Bézout's identity, probabilistic and
deterministic primality tests (culminating in the AKS theorem), integer
factorization algorithms (Pollard's rho, quadratic sieve, GNFS), discrete
logarithm algorithms (BSGS, Pohlig--Hellman, index calculus), and the
Chinese Remainder Theorem together with quadratic residues and continued
fractions. We carefully distinguish between \emph{proved theorems} (AKS,
Gauss reciprocity, Hasse's theorem) and \emph{open conjectures} (hardness
of factoring, DLP, RSA security). Cryptographic applications---RSA, Diffie--Hellman,
elliptic-curve cryptography, and post-quantum systems---are discussed in
terms of their mathematical foundations. All algorithms are implemented in
the companion Python module \texttt{algorithmic\_number\_theory.py} of the
Specialist Mathematics Project (Build 137).
\end{abstract}

\maketitle
\tableofcontents

% ===========================================================================
\section{Introduction}
% ===========================================================================

Number theory is one of the oldest branches of mathematics, yet it has become
central to modern computer science through the lens of algorithm design and
complexity theory. \emph{Algorithmic number theory} studies problems about
integers and their computational aspects: How quickly can we determine whether
a number is prime? How hard is it to factor a large integer? Can we compute
discrete logarithms efficiently?

These questions are not merely academic. They form the mathematical bedrock of
modern cryptography. RSA encryption, Diffie--Hellman key exchange, and
elliptic-curve cryptosystems all derive their security from the presumed
computational difficulty of certain number-theoretic problems. Understanding
both the algorithms and their complexity is therefore essential for
cryptographic design and analysis.

\subsection*{Scope of this paper}

We survey the principal algorithms of algorithmic number theory organized
around five core topics:
\begin{itemize}[leftmargin=2em]
  \item Greatest common divisors and the Extended Euclidean Algorithm
  \item Primality testing: from probabilistic to deterministic polynomial-time
  \item Integer factorization and its connection to cryptographic hardness
  \item Discrete logarithm algorithms
  \item Chinese Remainder Theorem, quadratic residues, and continued fractions
\end{itemize}
Each section presents the mathematical theory, the principal algorithms with
their complexity, and connections to cryptographic applications. We carefully
distinguish between results that are \emph{proved theorems} and statements
that are \emph{open conjectures}.

\subsection*{Notation}

Throughout, $\log$ denotes the natural logarithm unless otherwise noted.
$\mathbb{Z}_n = \mathbb{Z}/n\mathbb{Z}$ is the ring of residues modulo $n$.
$(\mathbb{Z}/n\mathbb{Z})^*$ denotes its multiplicative group, of order
$\varphi(n)$ (Euler's totient). Big-$O$ notation is standard;
$L_n[\alpha, c] = \exp\!\bigl(c(\log n)^\alpha (\log\log n)^{1-\alpha}\bigr)$
denotes subexponential complexity.

% ===========================================================================
\section{Greatest Common Divisor and the Extended Euclidean Algorithm}
% ===========================================================================

\subsection{The Euclidean Algorithm}

The \emph{greatest common divisor} (GCD) of two non-negative integers $a$ and
$b$, written $\gcd(a,b)$, is the largest integer dividing both. The Euclidean
algorithm computes $\gcd(a,b)$ via repeated division:
\begin{align*}
  a &= q_0 b + r_0,\quad 0 \le r_0 < b,\\
  b &= q_1 r_0 + r_1,\quad 0 \le r_1 < r_0,\\
  &\vdots\\
  r_{k-2} &= q_k r_{k-1} + 0.
\end{align*}
Then $\gcd(a,b) = r_{k-1}$, the last non-zero remainder.

\begin{theorem}[Euclidean Algorithm, Complexity]
The Euclidean algorithm terminates in $O(\log \min(a,b))$ division steps.
Each step reduces the pair $(a,b)$ to $(b, a \bmod b)$, and
$a \bmod b < a/2$ whenever $b \le a/2$, guaranteeing exponential convergence.
\end{theorem}

\begin{proof}
At each step, $r_{i+1} < r_i$. We claim $r_{i+2} < r_i/2$: if
$r_{i+1} \le r_i/2$ we are done; otherwise $r_{i+2} = r_i - r_{i+1} < r_i/2$.
Hence every two steps the remainder at least halves, so after
$2\lfloor\log_2 \min(a,b)\rfloor$ steps the algorithm terminates.
\end{proof}

\subsection{Extended Euclidean Algorithm and Bézout's Identity}

\begin{theorem}[Bézout's Identity]
For any integers $a,b$ (not both zero) there exist integers $s,t$ such that
\[
  sa + tb = \gcd(a,b).
\]
\end{theorem}

The \emph{Extended Euclidean Algorithm} computes $\gcd(a,b)$ and the
Bézout coefficients $s,t$ simultaneously by back-substituting through the
sequence of remainders. It runs in the same $O(\log\min(a,b))$ time.

\subsubsection*{Modular inverse}

A direct application: the modular inverse of $a$ modulo $m$ (i.e.\
$x$ with $ax \equiv 1 \pmod{m}$) exists if and only if $\gcd(a,m)=1$,
and equals the Bézout coefficient $s$ reduced modulo $m$. This is used
pervasively in RSA decryption and CRT reconstruction.

\begin{example}
Compute $7^{-1} \pmod{11}$: extended Euclidean gives $1 = 8\cdot 7 - 5\cdot 11$,
so $7^{-1} \equiv 8 \pmod{11}$. Verification: $7 \cdot 8 = 56 = 5\cdot 11 + 1 \equiv 1$.
\end{example}

\subsection{Binary GCD and Lehmer's Algorithm}

For multi-precision arithmetic (integers with $k$ machine words), schoolbook
division costs $O(k^2)$ operations. \emph{Lehmer's algorithm} reduces this to
$O(k\log^2 k\log\log k)$ by replacing full divisions with single-word
approximations. The \emph{binary GCD} avoids divisions entirely using shifts and
subtractions, which is efficient for hardware implementations.

% ===========================================================================
\section{Primality Testing}
% ===========================================================================

\subsection{Trial Division}

The simplest primality test: check all integers from $2$ to $\lfloor\sqrt{n}\rfloor$.
If none divide $n$, then $n$ is prime. Time: $O(\sqrt{n})$, exponential in the
number of digits. Practical only for $n \lesssim 10^{12}$.

\subsection{Fermat's Test and Pseudoprimes}

\begin{theorem}[Fermat's Little Theorem]
If $p$ is prime and $\gcd(a,p)=1$, then $a^{p-1} \equiv 1 \pmod{p}$.
\end{theorem}

Checking $a^{n-1} \equiv 1 \pmod{n}$ for several bases $a$ provides a fast
primality heuristic. However, \emph{Carmichael numbers} (e.g.\ $561 = 3\cdot11\cdot17$)
satisfy $a^{n-1}\equiv 1\pmod{n}$ for all $\gcd(a,n)=1$, so Fermat's test
can never certify them composite. There are infinitely many Carmichael numbers
(proved by Alford, Granville, and Pomerance in 1994).

\subsection{The Miller--Rabin Test}

\begin{definition}[Strong Pseudoprime]
Write $n-1 = 2^s \cdot d$ with $d$ odd. Then $n$ is a \emph{strong pseudoprime}
to base $a$ if either $a^d \equiv 1 \pmod{n}$ or $a^{2^r d} \equiv -1 \pmod{n}$
for some $0 \le r < s$.
\end{definition}

\begin{theorem}[Miller--Rabin, 1976]
If $n$ is composite, then the fraction of bases $a \in \{1,\ldots,n-1\}$ for
which $n$ is a strong pseudoprime to base $a$ is at most $1/4$. Consequently,
after $k$ independent random bases the probability of erroneously declaring
a composite as prime is at most $4^{-k}$.
\end{theorem}

With $k=40$ rounds the error probability falls below $10^{-24}$, sufficient for
all practical purposes. The deterministic variant (using specific bases depending
on the size of $n$) is unconditionally correct for all $n < 3.3 \times 10^{24}$.

\subsection{The Solovay--Strassen Test}

\begin{theorem}[Euler's Criterion]
For an odd prime $p$ and $\gcd(a,p)=1$:
$a^{(p-1)/2} \equiv \left(\frac{a}{p}\right) \pmod{p}$,
where $\bigl(\frac{a}{p}\bigr)$ is the Legendre symbol.
\end{theorem}

The \emph{Solovay--Strassen test} extends this using the Jacobi symbol:
check $a^{(n-1)/2} \equiv \bigl(\frac{a}{n}\bigr)_{\!J} \pmod{n}$. Unlike
Carmichael numbers with Fermat, there are no analogous universal
failures here: for composite $n$, at most half the bases pass.

\subsection{The AKS Primality Test}

\begin{theorem}[Agrawal--Kayal--Saxena, 2002]
\textbf{(Proved.)} \emph{PRIMES is in P.} There exists a deterministic algorithm
that decides primality of $n$ in time $\tilde{O}(\log^6 n)$.
\end{theorem}

The AKS test is based on the polynomial identity: $n$ is prime if and only if
$(X+a)^n \equiv X^n + a \pmod{X^r - 1, n}$ holds for all $a \le \lfloor\sqrt{\varphi(r)}\log n\rfloor$,
where $r$ is chosen so that the multiplicative order of $n$ modulo $r$ exceeds
$\log^2 n$. The algorithm proceeds as follows:

\begin{enumerate}[label=\arabic*.]
  \item \textbf{Perfect power test:} If $n = a^b$ for $b\ge 2$, return
    \textsc{Composite}.
  \item \textbf{Find $r$:} Find the smallest $r$ such that
    $\mathrm{ord}_r(n) > \log^2 n$.
  \item \textbf{Small factor test:} If $1 < \gcd(a,n) < n$ for some
    $a \le r$, return \textsc{Composite}.
  \item \textbf{Size check:} If $n \le r$, return \textsc{Prime}.
  \item \textbf{Polynomial test:} For $a = 1,\ldots, \lfloor\sqrt{\varphi(r)}\log n\rfloor$,
    check $(X+a)^n \equiv X^n + a \pmod{X^r-1, n}$.
  \item Return \textsc{Prime}.
\end{enumerate}

The correctness proof uses the theory of introspective numbers and a careful
counting argument. Although AKS established the theoretical landmark of
\textsc{PRIMES} $\in$ P, Miller--Rabin remains faster in practice.

\begin{remark}
The original complexity of $O(\log^{12} n)$ was improved to $O(\log^{6} n)$
by Lenstra and Pomerance (2005). Assuming the Generalized Riemann Hypothesis,
an earlier result of Miller already gave a polynomial-time deterministic test.
\end{remark}

% ===========================================================================
\section{Integer Factorization}
% ===========================================================================

\subsection{The Factorization Problem}

\begin{conjecture}[Hardness of Factoring]
\textbf{(Open.)} No polynomial-time algorithm exists for integer factorization.
Equivalently, the Integer Factorization Problem (IFP) is not in P (or even
in BPP). The connection to P~$\neq$~NP is subtle: IFP is in NP $\cap$ co-NP,
so it would not be NP-complete unless NP $=$ co-NP; but it might still be
strictly harder than P.
\end{conjecture}

This conjecture underpins RSA security. The best known classical algorithm
(GNFS, see below) runs in subexponential time $L_n[1/3, c]$; no polynomial
algorithm is known.

\subsection{Pollard's $\rho$ Algorithm}

\emph{Pollard's $\rho$ algorithm} (1975) finds a non-trivial factor of $n$
using a pseudo-random sequence $x_{i+1} = f(x_i) \bmod n$ (typically
$f(x) = x^2 + c$). By the birthday paradox, a collision $\gcd(|x_i - x_j|, n) > 1$
is expected after $O(n^{1/4})$ steps. Floyd's cycle detection reduces space
to $O(1)$.

\begin{theorem}[Expected Complexity of Pollard $\rho$]
Under heuristic assumptions on the pseudo-randomness of $f$, the expected
number of steps to find a factor $p \mid n$ is $O(p^{1/2}) = O(n^{1/4})$.
\end{theorem}

\subsection{Fermat Factoring and the Quadratic Sieve}

\emph{Fermat's method} exploits $n = a^2 - b^2 = (a-b)(a+b)$: search for
$a \ge \lceil\sqrt{n}\rceil$ with $a^2 - n$ a perfect square. Efficient when
$n$ has two factors of nearly equal size.

The \emph{quadratic sieve} (Pomerance, 1981) is the fastest algorithm for
$n$ up to about $10^{100}$. It seeks pairs $(x, y)$ with $x^2 \equiv y^2
\pmod{n}$, then $\gcd(x - y, n)$ reveals a factor. The sieving step finds
many $B$-smooth values of $x^2 \bmod n$ simultaneously, achieving heuristic
complexity $L_n[1/2, 1]$.

\subsection{The General Number Field Sieve}

The \emph{General Number Field Sieve} (GNFS) is currently the fastest known
algorithm for integers $n > 10^{100}$. It operates in an algebraic number
field $\mathbb{Z}[\theta]$ for a root $\theta$ of a carefully chosen
polynomial, sieving for smooth norms.

\begin{theorem}[GNFS Complexity (Heuristic)]
GNFS factors $n$ in expected time
\[
  L_n\!\left[\tfrac{1}{3},\, \left(\tfrac{64}{9}\right)^{1/3}\right]
  \approx \exp\!\left(1.923\,(\log n)^{1/3}(\log\log n)^{2/3}\right).
\]
\end{theorem}

The $512$-bit RSA-155 record ($1999$) required $8000$ MIPS-years; the
$829$-bit RSA-250 record ($2020$) required roughly $2700$ core-years.

\subsection{Implications for RSA}

\begin{conjecture}[RSA Hardness]
\textbf{(Open.)} Given a product $n = pq$ of two large primes and a ciphertext
$c = m^e \bmod n$, recovering $m$ without knowing $p,q$ is computationally
infeasible. No polynomial-time algorithm for the RSA problem is known. The
RSA problem is polynomial-time equivalent to factoring $n$ under natural
reductions.
\end{conjecture}

Current recommendations: RSA key lengths of $\ge 3072$ bits (NIST 2022)
or transition to elliptic-curve systems.

% ===========================================================================
\section{The Discrete Logarithm Problem}
% ===========================================================================

\subsection{Definition and Hardness}

\begin{definition}[Discrete Logarithm]
Let $g$ be a generator of the cyclic group $(\mathbb{Z}/p\mathbb{Z})^*$ of
order $p-1$. The \emph{discrete logarithm} of $h$ to base $g$ is the unique
$x \in \{0, 1, \ldots, p-2\}$ with $g^x \equiv h \pmod{p}$.
\end{definition}

\begin{conjecture}[DLP Hardness]
\textbf{(Open.)} No polynomial-time algorithm for the Discrete Logarithm
Problem (DLP) in $(\mathbb{Z}/p\mathbb{Z})^*$ is known. The best algorithms
achieve subexponential complexity $L_p[1/3, c]$, the same as GNFS.
\end{conjecture}

\subsection{Baby-Step Giant-Step (BSGS)}

\begin{theorem}[Shanks, 1971]
The Baby-Step Giant-Step algorithm computes $\log_g h$ in $(\mathbb{Z}/p\mathbb{Z})^*$
in $O(\sqrt{p})$ time and space.
\end{theorem}

Set $m = \lceil\sqrt{p-1}\rceil$ and write $x = im - j$:
\begin{itemize}
  \item \textbf{Baby steps:} Compute the table $\{(h \cdot g^j \bmod p, j) : 0 \le j < m\}$.
  \item \textbf{Giant steps:} Compute $g^{im} \bmod p$ for $i = 1, \ldots, m$.
  \item \textbf{Match:} When $g^{im} = h \cdot g^j$, we have $x = im - j$.
\end{itemize}

\subsection{Pohlig--Hellman Algorithm}

\begin{theorem}[Pohlig--Hellman, 1978]
If $p-1 = \prod_i q_i^{e_i}$ with all $q_i$ small, the discrete logarithm
in $(\mathbb{Z}/p\mathbb{Z})^*$ can be computed in
$O\bigl(\sum_i e_i(\log(p-1) + \sqrt{q_i})\bigr)$ operations.
\end{theorem}

The algorithm solves DLP modulo each prime power $q_i^{e_i}$ separately using
BSGS in the corresponding subgroup, then combines via the Chinese Remainder
Theorem. This makes DLP easy when $p-1$ is \emph{smooth} (all factors small).

\begin{remark}
This motivates the use of \emph{safe primes} $p = 2q+1$ (with $q$ also prime)
in Diffie--Hellman: then $p-1 = 2q$ has the large prime factor $q$, making
Pohlig--Hellman inapplicable.
\end{remark}

\subsection{Index Calculus}

The \emph{index calculus} algorithm achieves subexponential complexity for DLP
in $\mathbb{F}_p^*$. It proceeds in two phases:

\begin{enumerate}[label=\arabic*.]
  \item \textbf{Relation phase:} Build a factor base $\mathcal{B} = \{p_1,\ldots,p_k\}$
    of small primes. For random exponents $s$, check if $g^s \bmod p$ is
    $B$-smooth. When it is, record $s \equiv \sum_i a_i \log_g p_i \pmod{p-1}$.
  \item \textbf{Individual logarithm:} For target $h$, find $r$ with
    $h \cdot g^r$ $B$-smooth, then compute $\log_g h$.
\end{enumerate}

\begin{theorem}[Index Calculus Complexity]
Heuristically, index calculus computes discrete logarithms in
$(\mathbb{Z}/p\mathbb{Z})^*$ in time $L_p[1/2, 1]$; with the number field
sieve variant, in time $L_p[1/3, c]$.
\end{theorem}

% ===========================================================================
\section{The Chinese Remainder Theorem}
% ===========================================================================

\subsection{Statement and Proof}

\begin{theorem}[Chinese Remainder Theorem (CRT)]
Let $m_1, m_2, \ldots, m_k$ be pairwise coprime positive integers with
$M = m_1 m_2 \cdots m_k$. For any integers $r_1, r_2, \ldots, r_k$, the
system
\[
  x \equiv r_i \pmod{m_i}, \quad i = 1, \ldots, k,
\]
has a unique solution $x \pmod{M}$.
\end{theorem}

\begin{proof}
Set $M_i = M/m_i$ and let $e_i = M_i \cdot (M_i^{-1} \bmod m_i)$. Note that
$e_i \equiv 1 \pmod{m_i}$ and $e_i \equiv 0 \pmod{m_j}$ for $j \ne i$.
Then $x = \sum_{i=1}^k r_i e_i \bmod M$ satisfies all congruences. Uniqueness
follows because any two solutions differ by a multiple of $M$.
\end{proof}

\subsection{Algorithmic Complexity}

The CRT reconstruction algorithm runs in $O(k \log^2 M)$ bit operations (using
fast modular arithmetic). For $k = O(1)$ modules, this reduces to
$O(\log^2 M)$, i.e.\ quasi-linear in the number of digits.

\subsection{Cryptographic Applications}

\subsubsection*{RSA-CRT decryption}

Instead of computing $c^d \bmod n$ with $d \approx n$, compute
$c^{d_p} \bmod p$ and $c^{d_q} \bmod q$ separately (where
$d_p = d \bmod (p-1)$, $d_q = d \bmod (q-1)$), then combine via CRT.
This gives a roughly $4\times$ speedup (Garner's formula).

\subsubsection*{Multi-modular arithmetic}

Large integer multiplication can be performed by working modulo several
machine-word-sized primes and reconstructing via CRT, avoiding multi-precision
arithmetic entirely.

% ===========================================================================
\section{Quadratic Residues and the Tonelli--Shanks Algorithm}
% ===========================================================================

\subsection{Legendre Symbol and Euler's Criterion}

\begin{definition}[Legendre Symbol]
For an odd prime $p$ and integer $a$ with $\gcd(a,p)=1$:
\[
  \left(\frac{a}{p}\right) =
  \begin{cases} 1 & \text{if } a \text{ is a quadratic residue mod } p,\\
                -1 & \text{otherwise.}
  \end{cases}
\]
\end{definition}

\begin{theorem}[Euler's Criterion]
$\left(\frac{a}{p}\right) \equiv a^{(p-1)/2} \pmod{p}$.
\end{theorem}

\subsection{The Gauss Quadratic Reciprocity Law}

\begin{theorem}[Gauss, 1796. \textbf{Proved.}]
For distinct odd primes $p, q$:
\[
  \left(\frac{p}{q}\right)\left(\frac{q}{p}\right) =
  (-1)^{\tfrac{p-1}{2}\cdot\tfrac{q-1}{2}}.
\]
Equivalently, $\left(\frac{p}{q}\right) = \left(\frac{q}{p}\right)$ unless
$p \equiv q \equiv 3 \pmod{4}$, in which case they have opposite signs.
\end{theorem}

Gauss gave eight proofs. Together with the first and second supplements
(concerning $(-1/p)$ and $(2/p)$), quadratic reciprocity enables efficient
computation of the Legendre symbol in $O(\log^2 p)$ steps.

\subsection{Jacobi Symbol}

The \emph{Jacobi symbol} $\bigl(\frac{a}{n}\bigr)_J$ extends the Legendre
symbol to composite $n$ via multiplicativity over the prime factorization of
$n$. It can be computed by reciprocity laws in $O(\log^2 n)$ without
factoring $n$. However, $\bigl(\frac{a}{n}\bigr)_J = 1$ does \emph{not}
imply $a$ is a QR mod $n$ when $n$ is composite.

\subsection{Tonelli--Shanks Algorithm}

Given a quadratic residue $a$ modulo an odd prime $p$, the
\emph{Tonelli--Shanks algorithm} computes $\sqrt{a} \bmod p$ (i.e.\ $x$ with
$x^2 \equiv a \pmod p$) in $O(\log^2 p)$ expected time.

Write $p-1 = 2^s \cdot q$ with $q$ odd. Find a non-residue $z$ and initialize
$M = s$, $c = z^q \bmod p$, $t = a^q \bmod p$, $R = a^{(q+1)/2} \bmod p$.
The algorithm iterates, halving $M$ and adjusting $R$ via $c$, until $t = 1$.

\begin{theorem}[Tonelli--Shanks Complexity]
The expected running time is $O(\log^2 p)$ modular multiplications, where
the expectation is over the random choice of the non-residue $z$.
\end{theorem}

% ===========================================================================
\section{Continued Fractions}
% ===========================================================================

\subsection{Definition and Periodic Expansion}

\begin{definition}[Continued Fraction]
A \emph{simple continued fraction} is an expression
$[a_0; a_1, a_2, \ldots] = a_0 + \cfrac{1}{a_1 + \cfrac{1}{a_2 + \cdots}}$
with $a_i \in \mathbb{Z}$, $a_i > 0$ for $i \ge 1$.
\end{definition}

\begin{theorem}[Lagrange, 1770. Proved.]
The continued fraction expansion of $\sqrt{n}$ (for $n$ not a perfect square)
is eventually periodic:
$\sqrt{n} = [a_0; \overline{a_1, a_2, \ldots, a_k}]$ with period length
$k \le 2a_0$.
\end{theorem}

The convergents $p_k/q_k$ of $\sqrt{n}$ satisfy $|p_k - q_k\sqrt{n}| < 1/q_{k+1}$,
making them the best rational approximations to $\sqrt{n}$.

\subsection{Pell's Equation}

\begin{theorem}[Lagrange, Proved]
The equation $x^2 - ny^2 = 1$ (Pell's equation, $n$ not a perfect square)
has infinitely many integer solutions. The fundamental solution $(x_1, y_1)$
is given by the convergents of $\sqrt{n}$: if the period length $k$ is even,
then $(p_{k-1}, q_{k-1})$ is the fundamental solution; if $k$ is odd, then
$(p_{2k-1}, q_{2k-1})$.
\end{theorem}

\begin{example}
$n=2$: $\sqrt{2} = [1;\overline{2}]$, period $k=1$. Fundamental solution:
$(p_1, q_1) = (3, 2)$: $3^2 - 2\cdot 2^2 = 9 - 8 = 1$.
\end{example}

\subsection{CFRAC Factorization}

The \emph{continued fraction factoring algorithm} (Morrison--Brillhart, 1975)
is an early smooth-number algorithm. It seeks pairs $(x,y)$ with
$x^2 \equiv y^2 \pmod{n}$ by collecting the numerators $p_k$ of the
convergents of $\sqrt{n}$ and sieving for $B$-smooth values of
$p_k^2 \bmod n$. This was among the first algorithms with
subexponential heuristic complexity, predating the quadratic sieve.

\subsection{Rational Approximations and Cryptanalysis}

Continued fractions appear in the Wiener attack on RSA: if the secret
exponent $d < n^{1/4}$, then $d$ can be recovered from the continued
fraction expansion of $e/n$.

% ===========================================================================
\section{Cryptographic Applications}
% ===========================================================================

\subsection{RSA Cryptosystem}

\begin{definition}[RSA]
Key generation: choose large primes $p,q$, set $n = pq$, $\varphi(n) = (p-1)(q-1)$.
Choose public exponent $e$ with $\gcd(e, \varphi(n)) = 1$; compute
$d = e^{-1} \bmod \varphi(n)$ (Extended Euclidean). \textbf{Public key:} $(n, e)$.
\textbf{Private key:} $d$. \textbf{Encryption:} $c = m^e \bmod n$.
\textbf{Decryption:} $m = c^d \bmod n$ (follows from $ed \equiv 1 \pmod{\varphi(n)}$
and Euler's theorem).
\end{definition}

RSA security rests on the hardness of factoring $n$ (and possibly further
assumptions). See Section~4.4.

\subsection{Diffie--Hellman Key Exchange}

\begin{definition}[DH Protocol, 1976]
Public: large prime $p$, generator $g$ of $(\mathbb{Z}/p\mathbb{Z})^*$.
Alice sends $g^a \bmod p$, Bob sends $g^b \bmod p$. Both compute the shared
secret $g^{ab} \bmod p$. An eavesdropper must solve the DLP to recover $a$ or $b$.
\end{definition}

\begin{conjecture}[Computational Diffie--Hellman, CDH]
\textbf{(Open.)} Given $g^a$ and $g^b$ in $(\mathbb{Z}/p\mathbb{Z})^*$,
computing $g^{ab}$ is hard. The CDH assumption is at least as hard as DLP.
\end{conjecture}

\subsection{Elliptic Curve Cryptography (ECC)}

Let $E: y^2 = x^3 + ax + b$ over $\mathbb{F}_p$. The group $E(\mathbb{F}_p)$
under the chord-and-tangent law has order approximately $p$ (Hasse's theorem).
The elliptic curve DLP (ECDLP) is believed harder than the classical DLP for
the same group size, as no subexponential general algorithm is known for
elliptic curves.

\begin{theorem}[Hasse, 1933. Proved.]
$|{|E(\mathbb{F}_p)| - (p+1)}| \le 2\sqrt{p}$.
\end{theorem}

ECDSA (Elliptic Curve Digital Signature Algorithm) and ECDH use 256-bit keys
for security equivalent to RSA-3072, a dramatic efficiency gain.

\subsection{Post-Quantum Cryptography}

Shor's quantum algorithm (1994) solves both integer factorization and DLP in
$O((\log n)^3)$ quantum time, breaking RSA, DH, and ECC. This motivates
\emph{post-quantum cryptography}:

\begin{itemize}
  \item \textbf{Lattice-based:} CRYSTALS-Kyber (NIST PQC standard, 2022)
  \item \textbf{Hash-based:} SPHINCS+
  \item \textbf{Code-based:} McEliece (1978)
  \item \textbf{Multivariate:} Rainbow (broken), HQC
\end{itemize}

NIST standardized the first post-quantum algorithms in 2022--2024. The
mathematical hardness of lattice problems (LWE, NTRU) rests on conjectured
worst-case to average-case reductions.

% ===========================================================================
\section{Conclusion and Connections}
% ===========================================================================

Algorithmic number theory sits at the intersection of pure mathematics and
theoretical computer science. The core results cluster around three themes:

\medskip
\noindent\textbf{Proved theorems} that underpin the algorithms:
Euclidean algorithm ($O(\log n)$ steps), Bézout's identity, Fermat's little
theorem, AKS primality in $P$, Gauss quadratic reciprocity, Lagrange on Pell's
equation, Hasse's theorem on elliptic curves.

\medskip
\noindent\textbf{Widely believed but unproved conjectures}:
The hardness of factoring (IFP $\notin$ P), DLP hardness (CDH), and
RSA security. These are the load-bearing assumptions of modern cryptography.
The connection to P~$\ne$~NP is indirect: IFP $\in$ NP $\cap$ co-NP, so
it would not be NP-complete unless the polynomial hierarchy collapses.

\medskip
\noindent\textbf{Open problems}:
Twin prime conjecture, Goldbach's conjecture, distribution of prime gaps
(Cramér's conjecture), efficient algorithms for ECDLP.

The transition to post-quantum cryptography, driven by Shor's algorithm,
illustrates how advances in algorithmic number theory directly reshape
the security landscape of global communications.

% ===========================================================================
\begin{thebibliography}{99}
% ===========================================================================

\bibitem{AKS2002}
M.~Agrawal, N.~Kayal, N.~Saxena,
\emph{PRIMES is in P},
Ann.\ of Math.\ \textbf{160} (2004), 781--793.

\bibitem{CP2005}
R.~Crandall, C.~Pomerance,
\emph{Prime Numbers: A Computational Perspective},
2nd ed., Springer, New York, 2005.

\bibitem{Cohen1993}
H.~Cohen,
\emph{A Course in Computational Algebraic Number Theory},
Graduate Texts in Mathematics \textbf{138}, Springer, 1993.

\bibitem{MenezesOV1996}
A.~Menezes, P.~van Oorschot, S.~Vanstone,
\emph{Handbook of Applied Cryptography},
CRC Press, 1996. (Available at \url{https://cacr.uwaterloo.ca/hac/}).

\bibitem{Shoup2009}
V.~Shoup,
\emph{A Computational Introduction to Number Theory and Algebra},
2nd ed., Cambridge University Press, 2009.

\bibitem{Pollard1975}
J.~M.~Pollard,
\emph{A Monte Carlo method for factorization},
BIT Numerical Mathematics \textbf{15} (1975), 331--334.

\bibitem{Miller1976}
G.~L.~Miller,
\emph{Riemann's hypothesis and tests for primality},
J.~Comput.\ System Sci.\ \textbf{13} (1976), 300--317.

\bibitem{PohligHellman1978}
S.~C.~Pohlig, M.~E.~Hellman,
\emph{An improved algorithm for computing logarithms over $GF(p)$ and its cryptographic significance},
IEEE Trans.\ Inform.\ Theory \textbf{24} (1978), 106--110.

\bibitem{Shanks1971}
D.~Shanks,
\emph{Class Number, a Theory of Factorization and Genera},
Proc.\ Symp.\ Pure Math.\ \textbf{20}, AMS, 1971, pp.~415--440.

\bibitem{Shor1994}
P.~W.~Shor,
\emph{Algorithms for quantum computation: discrete logarithms and factoring},
Proc.\ 35th FOCS, 1994, pp.~124--134.

\end{thebibliography}

\end{document}
